\section{Future Work}
\label{sec:future}

P0, P1, P2 and the bounded Subject-C cross-repository replication step are complete and must not be retrospectively rerun to obtain cleaner outcomes. Future work should extend the evidence population rather than modify accepted samples.

The highest-priority next scientific study is a larger preregistered distributional experiment with more independent repetitions, additional repositories and task families, and a dependence-aware statistical analysis fixed before outcomes. If equivalence or non-inferiority is a target, the margin must be justified in engineering terms and preregistered. The current P1/P2/Subject-C results cannot supply that margin post hoc.

The programme's full Level-3 external-validity threshold remains future work. One Subject-C cross-repository replication does not provide variation across multiple unrelated repositories, languages, architectures, repository sizes, dependency widths and documentation/test quality. Independent task selection and external replication would further reduce author/repository bias.

Reconstruction economics should be tested directly. A future study should measure reconstruction-attributed tokens, files/bytes read, searches, model/tool calls, elapsed time, retries and actual metered charges where exposed. It must count rediscovery as work and must not infer provider-internal GPU/KV-cache economics from client-side counters.

Semantic-state ablation should remove small preregistered information classes rather than broadly degrading the repository. The scientific question is which durable artefacts have marginal continuity value.

Tiered execution belongs in a separate experiment because changing executor capability would confound the core history treatment. Security claims likewise require adversarial evaluation rather than architectural argument alone.

Cross-model continuity remains particularly interesting. A design in which model family $M_A$ produces an accepted state and a different family $M_B$ resumes from it without predecessor model-native state would provide stronger evidence that continuity resides in external engineering state rather than a provider/model-specific session representation.

Reconstruction scaling should vary repository mass and dependency width while holding task semantics as constant as possible. Plain search, symbol indexes, static-analysis graphs, repository maps and model-guided retrieval can be compared without claiming any one is intrinsically RaS.

The broader Authoritative-State Externalisation Principle remains speculative. Software repositories have unusually strong versioning, diffability and verification properties; other domains may depend more heavily on tacit, continuously changing or externally governed information.
